\section{Klasični kriptosistemi, LPT}
\textbf{Def.}
\emph{Kriptosistem} je peterka \((\B, \CC, \K, \E, \D)\), kjer
\begin{itemize}
    \item \(\B\) je končna množica besedil;
    \item \(\CC\) je množica kriptogramov (angl.\ ciphertext);
    \item \(\K\) je množica ključev;
    \item \(\E = \setb{E_k: \B \to \CC}{k \in \K}\) je množica \emph{šifrirnih} polinomskih funkcij;
    \item \(\D = \setb{D_k: \CC \to \B}{k \in \K}\) je množica \emph{dešifrirnih} polinom.\ funkcij.
\end{itemize}
Pri tem za kriptosistem mora veljati \emph{pravilnost}, tj.\
\[
    \all{k \in \K} \some{k' \in \K} \all{m \in \B} D_{k'}(E_k(m)) = m.
\]
%
\s{Popolna tajnost}\\
\textbf{Def.} Kriptosistem \emph{ima lastnost popolne tajnosti (LPT)}, če 
    \[
        \all{m \in \B} \all{c \in \CC} P(X_\B = m\, |\, X_\CC = c) = P(X_\B = m).
    \]
\textbf{Lema.} Kriptosistem ima LPT \(\liff P(X_\CC = c\, |\, X_\B = m) = P(X_\CC = c)\).\\
\textbf{Trd.} Če ima kriptosistem LPT, potem \(|\B| \leq |\CC| \leq |\K|\).\\
\begin{itemize}
    \item \(P(X_\CC = c) = \sum_{b \in \B}P(X_\CC = c \,|\, X_\B = b) \cdot P(X_\B = b)\)
\end{itemize}
\s{Vernamova šifra (OTP)}\\
Naj bodo \(\B = \CC = \K = \set{0, 1}^\lambda,\ \lambda > 0\). Ključi izbiramo enakomerno naključno. Definiramo 
\begin{itemize}
    \item \(E_k(m) = m \oplus k\); 
    \item \(D_k(c) = c \oplus k\).
\end{itemize}
\s{Afina šifra}\\
Naj bodo \(\B = \CC = \Z_m\) in \(\K = \Z_m^* \times \Z_m\). Definiramo
\begin{itemize}
    \item \(E_{(a, b)}(x) = (ax + b) \mod m\);
    \item \(D_{(a, b)}(y) = a^{-1}(y - b) \mod m\).
\end{itemize}
\s{Hillova šifra}\\
Naj bodo \(\B = \CC = \Z_m^n\) in \(\K = \setb{A \in \Z_m^{n \times n}}{\det A \in \Z_m^*}\). Definiramo
\begin{itemize}
    \item \(E_K(B) = K \cdot B \mod m\);
    \item \(D_K(C) = K^{-1} \cdot C \mod m\).
\end{itemize}
Če ne poznamo dolžino bloka jo poskusimo uganiti \((2, 3, \ldots)\). Za razbitje v splošnem potrebujemo \(n\) parov (besedilo/kriptogram).
